% ============================================================================
%  Rozdział 1. Wstęp
% ============================================================================
\chapter{Wstęp}

Gry zręcznościowe należą do najstarszych gatunków gier komputerowych, a~mimo to wciąż
stanowią dużą część rynku, zwłaszcza na urządzeniach mobilnych. O~ich odbiorze decyduje
przede wszystkim to, czy sterowanie jest natychmiastowe i~przewidywalne, czy reguły można
zrozumieć bez instrukcji oraz czy pojedyncza rozgrywka trwa na tyle krótko, aby gracz chciał
ją powtórzyć. Spełnienie tych warunków wymaga dopracowania mechaniki ruchu, reakcji na
kolizje i~informacji zwrotnej w~interfejsie, czyli elementów, które w~kodzie gry są ze sobą
silnie powiązane.

Wraz ze wzrostem liczby obiektów i~zdarzeń na planszy rośnie liczba zależności pomiędzy
fragmentami kodu. Jeśli każdy skrypt sam wyszukuje potrzebne mu obiekty i~bezpośrednio
wywołuje metody innych komponentów, już przy kilkudziesięciu klasach zmiana jednej
mechaniki zaczyna wymagać poprawek w~wielu miejscach. Niniejsza praca opisuje projekt
i~implementację trójwymiarowej gry zręcznościowej, w~której od początku przyjęto inne podejście:
zależności są dostarczane przez kontener wstrzykiwania zależności, zdarzenia rozgrywki są
modelowane jako strumienie reaktywne, a~logika gry jest oddzielona od komponentów silnika
Unity.

\section{Uzasadnienie podjęcia tematu pracy}

Silnik Unity jest jednym z~najpowszechniej stosowanych narzędzi do tworzenia gier
niezależnych i~mobilnych. Jego model komponentowy pozwala szybko zbudować działający
prototyp, jednak nie narzuca żadnej architektury kodu. Typowy projekt rozwijany bez
przyjętych wcześniej zasad kończy się dużą liczbą klas dziedziczących po
\klasa{MonoBehaviour}, które odwołują się do siebie przez \kod{FindObjectOfType}, pola
publiczne ustawiane w~inspektorze lub statyczne singletony. Taki kod trudno testować,
a~kolejność inicjalizacji obiektów zależy od kolejności wywołań metod cyklu życia, którą
łatwo przeoczyć \cite{unity-order}.

Wybór gry zręcznościowej jako przedmiotu pracy wynikał z~tego, że gatunek ten dobrze
eksponuje wspomniane problemy. Pojedyncza kolizja kuli z~monetą powinna jednocześnie
uruchomić animację znikania monety, zwiększyć licznik w~interfejsie, zaktualizować statystyki
poziomu oraz, jeśli była to ostatnia moneta, odblokować metę. Zdarzeń tego rodzaju jest
w~grze kilkadziesiąt i~zachodzą one w~różnych momentach klatki. Podejście reaktywne,
w~którym zdarzenie publikowane jest w~jednym miejscu, a~zainteresowane moduły
subskrybują je niezależnie od siebie, pozwala opisać takie zależności w~sposób jawny.

Dodatkowym powodem było zainteresowanie autora praktykami stosowanymi w~studiach
tworzących gry mobilne, w~których wstrzykiwanie zależności (Zenject), biblioteki reaktywne
(UniRx, a~obecnie R3) i~asynchroniczność oparta na \kod{async}/\kod{await} (UniTask) są
standardowym wyposażeniem projektu. Praca była okazją do sprawdzenia, jak te narzędzia
sprawdzają się w~kompletnej, choć niewielkiej grze, oraz jakie koszty projektowe
wprowadzają.

\section{Cel pracy}

Celem pracy było zaprojektowanie i~zaimplementowanie trójwymiarowej gry zręcznościowej
w~silniku Unity, w~której gracz steruje kulą poruszającą się po planszy ograniczonej
ścianami. Zadaniem gracza jest zebranie wszystkich monet znajdujących się na poziomie,
po czym pojawia się meta umożliwiająca ukończenie poziomu. Na trasie rozmieszczono
przeciwników i~przeszkody; zderzenie z~przeciwnikiem odbiera punkty życia, a~ich utrata
kończy rozgrywkę. Cel główny podzielono na następujące cele szczegółowe:
\begin{itemize}[leftmargin=1.9cm, itemsep=1pt]
  \item opracowanie koncepcji gry, zasad ukończenia poziomu i~zestawu mechanik,
  \item zaprojektowanie modułowej architektury aplikacji z~wykorzystaniem
        wstrzykiwania zależności (biblioteka Zenject) i~podziału kodu na niezależne zestawy,
  \item implementacja ruchu kuli, kolizji i~odbić z~wykorzystaniem systemu fizyki silnika,
  \item implementacja przeciwników o~zachowaniach autonomicznych oraz systemu punktów
        życia i~leczenia,
  \item przygotowanie interfejsu użytkownika w~podejściu MVVM, w~którym widok jest
        aktualizowany przez subskrypcje strumieni reaktywnych,
  \item ograniczenie alokacji pamięci w~trakcie rozgrywki przez wielokrotne użycie obiektów
        efektów cząsteczkowych (wzorzec \emph{Object Pool}),
  \item zapis wyników rozgrywki oraz integracja z~usługami zewnętrznymi (Firebase,
        Appodeal) na poziomie inicjalizacji,
  \item przygotowanie wersji instalacyjnej na urządzenia z~systemem Android i~przetestowanie
        działania gry.
\end{itemize}

Założono, że gra ma być kompletna z~punktu widzenia gracza (menu, wybór poziomu, ekran
ładowania, interfejs rozgrywki, okno zakończenia i~statystyki), natomiast nacisk w~części
opisowej zostanie położony na rozwiązania architektoniczne i~sposób ich realizacji w~kodzie.

\section{Zakres pracy}

Zakres pracy obejmował zarówno część koncepcyjną, jak i~wykonanie działającej aplikacji.
Prace rozpoczęto od określenia mechaniki sterowania. Zrezygnowano ze sterowania klawiaturą
na rzecz gestu przeciągnięcia od kuli: długość przeciągnięcia wyznacza siłę, a~jego
kierunek, odwrócony o~180°, kierunek ruchu. Takie sterowanie działa w~ten sam sposób
z~myszą i~na ekranie dotykowym, co ułatwiło późniejsze przeniesienie gry na system Android.

Następnie przygotowano szkielet architektury: kontekst projektu i~konteksty scen biblioteki
Zenject, instalatory dla poszczególnych obszarów gry, mechanizm wczytywania konfiguracji
z~plików JSON oraz podział kodu na zestawy kompilacji (\emph{assembly definitions}).
Na tej podstawie zaimplementowano kolejne mechaniki:
\begin{itemize}[leftmargin=1.9cm, itemsep=1pt]
  \item ruch kuli, tarcie, odbicia od ścian i~przeciwników wraz z~animacją zgniecenia kuli,
  \item monety, fabrykę monet i~warunek odblokowania mety,
  \item dwa typy przeciwników: patrolujące kaktusy i~nieruchome grzyby,
  \item punkty życia, apteczki leczące kulę, okno przegranej i~wygranej,
  \item pola przyspieszające, odbijacze, obracające się belki i~ruchome ściany,
  \item losowy wiatr z~wskaźnikiem kierunku, minimapę i~przycisk wypychający zablokowaną kulę,
  \item pomiar czasu przejścia poziomu, zapis najlepszego czasu i~okno statystyk w~menu.
\end{itemize}

Przygotowano trzy poziomy o~rosnącej trudności, scenę menu głównego, scenę ładowania oraz
scenę rozruchową inicjalizującą usługi zewnętrzne. Ostatnim etapem było przygotowanie
procesu budowania wersji na system Android, testy działania oraz opracowanie niniejszej
dokumentacji. Rysunek~\ref{rys:hud-wstep} przedstawia widok rozgrywki na pierwszym poziomie.

\begin{figure}[!htb]
  \centering
  \zrzut[0.72\textwidth]{poziom1_hud}
  \caption{Widok rozgrywki na pierwszym poziomie opracowanej gry; \zrodloWlasne}
  \label{rys:hud-wstep}
\end{figure}

\section{Problematyka pracy}

Pierwszym obszarem problemowym było pogodzenie fizyki silnika z~oczekiwaniami gracza
wobec gry zręcznościowej. Symulacja ciała sztywnego w~silniku PhysX traktuje kulę jak
obiekt fizyczny: prędkość maleje w~sposób zależny od materiałów, a~odbicia zależą od
współczynników sprężystości i~kąta zderzenia. Gracz oczekuje natomiast, że kula po
uderzeniu w~ścianę odbije się zgodnie z~zasadą „kąt padania równa się kątowi odbicia”
i~straci przewidywalną część prędkości. Z~tego powodu ruch kuli jest wyznaczany przez
własny model (kierunek i~moc pchnięcia), a~silnik fizyki odpowiada głównie za wykrywanie
kolizji i~rozwiązywanie przenikania obiektów.

Drugim problemem było sterowanie na ekranie dotykowym. Gest musi rozpoczynać się na kuli,
a~nie w~dowolnym miejscu ekranu, ponieważ w~przeciwnym razie przypadkowe dotknięcie
wystrzeliwałoby kulę. Wymagało to rzutowania promienia z~kamery na scenę i~pominięcia
obiektów typu \emph{trigger} (pól przyspieszających, apteczek, mety), które mogłyby
przesłonić kulę stojącą wewnątrz nich. Dodatkowo kula zatrzymana za ścianą bywa
niewidoczna dla gracza, co wymusiło dodanie przycisku awaryjnego wypchnięcia.

Trzecim obszarem była organizacja przepływu zdarzeń. W~grze występują łańcuchy zdarzeń
obejmujące kilka modułów (kolizja $\rightarrow$ animacja $\rightarrow$ licznik $\rightarrow$
meta $\rightarrow$ okno końca poziomu). Każda subskrypcja musi zostać zwolniona wraz ze
zniszczeniem obiektu, inaczej dochodzi do wycieków pamięci i~wywołań na nieistniejących
obiektach po przeładowaniu sceny. Konieczne było więc konsekwentne zarządzanie czasem
życia subskrypcji przez interfejs \klasa{IDisposable}.

Czwarty problem wynikał z~samego silnika: obiekty \klasa{MonoBehaviour} są tworzone przez
Unity, więc nie można przekazać im zależności przez konstruktor, a~obiekty tworzone
w~trakcie gry muszą powstawać przez kontener zależności, co wpłynęło na projekt fabryk.
Ostatnim zagadnieniem była zgodność z~kompilacją IL2CPP stosowaną na Androidzie, która
usuwa nieużywany kod i~może zablokować deserializację plików konfiguracyjnych wykonywaną
przez refleksję.

Oceniając realizację projektu, należy wskazać różnice pomiędzy wnioskiem o~temat a~wykonaną
aplikacją. Wniosek przewidywał grę desktopową i~bibliotekę UniRx. W~trakcie prac platformą
docelową stał się Android, ponieważ sterowanie gestem okazało się naturalne dla ekranu
dotykowego, a~zamiast UniRx zastosowano bibliotekę R3 tego samego autora, rozwijaną jako
jej następca \cite{r3,unirx}. Wersja uruchomiona w~edytorze Unity nadal jest w~pełni grywalna
przy użyciu myszy. Integracja z~Firebase została ograniczona do inicjalizacji i~weryfikacji
zależności, a~zapis danych obejmuje wyniki poziomów przechowywane lokalnie. Szczegółowe
omówienie tych decyzji zawiera rozdział~\ref{rozdz:zakonczenie}.

\section{Przegląd rynku i stan dostępnych rozwiązań}

Gry, w~których gracz kieruje toczącą się kulą przez plansze z~przeszkodami, są obecne na
rynku od dziesięcioleci. Różnią się przede wszystkim sposobem sterowania: część z~nich
pozwala bezpośrednio kierować kulą, inne każą przechylać planszę lub nadawać kuli
jednorazowy impuls, jak w~grach golfowych. W~tej części omówiono cztery tytuły, które pod
różnymi względami są bliskie projektowi opisanemu w~pracy.

\subsection{Super Monkey Ball}

\emph{Super Monkey Ball} to seria gier wydawana przez firmę Sega, zapoczątkowana w~2001
roku grą opracowaną przez studio Amusement Vision, wydaną najpierw na automaty, a~następnie
na konsolę GameCube \cite{gra-smb}. Gracz nie steruje bezpośrednio postacią zamkniętą
w~przezroczystej kuli, lecz przechyla całą planszę, przez co kula toczy się pod wpływem
grawitacji. Celem jest dotarcie do mety, zanim upłynie czas przeznaczony na planszę;
upływ czasu lub spadek z~planszy kosztuje życie, a~zbierane po drodze banany zwiększają
wynik.

Sterowanie przechyłem planszy wymaga precyzyjnego kontrolera analogowego; jego
odpowiednik na telefonie, oparty na czujnikach ruchu, był sprawdzany również w~niniejszym
projekcie, ale został odrzucony po testach (podrozdział~\ref{sec:poprawki}).

\subsection{Marble It Up!}

\emph{Marble It Up!} to gra zespołu The Marble Collective wydana w~2018 roku na konsolę
Nintendo Switch i~komputery osobiste, rozwinięta w~2023 roku do wersji \emph{Marble It Up!
Ultra} \cite{gra-miu}. Gracz bezpośrednio steruje toczącą się kulką na trójwymiarowych
planszach, z~których można spaść. Podstawą rywalizacji jest czas przejścia poziomu: za
szybkie ukończenie przyznawane są medale, wyniki trafiają do globalnych tablic
rankingowych, a~gracz może ścigać się z~zapisem wcześniejszego przejazdu.

Ten sam mechanizm powrotu do ukończonych poziomów w~celu poprawy wyniku zastosowano
w~opisywanej grze w~postaci licznika czasu i~zapisu najlepszego czasu każdego poziomu.

\subsection{Golf Battle}

\emph{Golf Battle} to mobilna gra w~mini-golfa wydana przez firmę Miniclip w~2018 roku na
systemy Android i~iOS \cite{gra-golf}. Strzał wykonuje się przez odciągnięcie palca od
piłki i~jego zwolnienie: długość przeciągnięcia określa siłę uderzenia, a~kierunek
przeciwny do ruchu palca wyznacza tor piłki \cite{gra-golf-recenzja}. Rozgrywka toczy się
w~trybie wieloosobowym, w~którym na tym samym torze może jednocześnie grać do sześciu
osób.

Sterowanie przyjęte w~opracowanej grze działa na tej samej zasadzie, ale strzał można
wykonać również w~czasie ruchu kuli, a~na planszy działają przeciwnicy, wiatr i~elementy
zmieniające jej prędkość.

\subsection{Angry Birds}

\emph{Angry Birds} fińskiego studia Rovio Entertainment, wydana w~grudniu 2009 roku na
system iOS, a~później na wiele innych platform, spopularyzowała na ekranach dotykowych
mechanikę procy \cite{gra-ab}. Gracz odciąga ptaka na gumie procy,
a~po zwolnieniu ptak leci po torze balistycznym. Siła i~kąt rzutu wynikają wprost
z~położenia palca względem procy, a~ślad poprzedniego lotu pozostaje na ekranie, co
pozwala korygować kolejne strzały.

Gra pokazała, że sterowanie oparte na odciągnięciu i~zwolnieniu jest zrozumiałe bez
samouczka; w~opisywanym projekcie kierunek strzału podpowiada linia celownicza.

\begin{table}[!htb]
  \centering
  \caption{Porównanie wybranych gier z~opracowaną grą; \zrodloWlasne}
  \label{tab:porownanie-gier}
  \begin{tabularx}{\textwidth}{>{\raggedright\arraybackslash}p{2.35cm}
      >{\raggedright\arraybackslash}X >{\raggedright\arraybackslash}X
      >{\raggedright\arraybackslash}X >{\centering\arraybackslash}p{1.45cm}}
    \toprule
    \textbf{Gra} & \textbf{Sterowanie} & \textbf{Cel poziomu} & \textbf{Zagrożenia} & \textbf{Punkty życia} \\
    \midrule
    Super Monkey Ball & przechylanie planszy & dotarcie do mety w~limicie czasu & spadek z~planszy, upływ czasu & nie \\
    Marble It Up! & bezpośrednie sterowanie kulką & dotarcie do mety, jak najkrótszy czas & spadek z~planszy & nie \\
    Golf Battle & przeciągnięcie i~zwolnienie & trafienie do dołka & przeszkody na torze & nie \\
    Angry Birds & odciągnięcie procy & zniszczenie celów & brak & nie \\
    \midrule
    Opracowana gra & przeciągnięcie i~zwolnienie, także w~ruchu & zebranie monet, dotarcie do mety & przeciwnicy, przeszkody ruchome, wiatr & tak \\
    \bottomrule
  \end{tabularx}
\end{table}

Zestawienie w~tabeli~\ref{tab:porownanie-gier} pokazuje, że opracowana gra łączy sterowanie
znane z~gier golfowych z~celem i~strukturą poziomów typowymi dla gier z~toczącą się kulą.
Elementem wyróżniającym jest system punktów życia: kontakt z~przeciwnikiem nie kończy od
razu rozgrywki, lecz zmniejsza zapas zdrowia, który można odnowić, zbierając apteczki.

\section{Układ i struktura pracy}

Praca została podzielona na sześć rozdziałów. Rozdział pierwszy ma charakter wprowadzający:
przedstawia uzasadnienie wyboru tematu, cel i~zakres pracy, problematykę oraz przegląd
podobnych gier. W~rozdziale drugim omówiono podstawy teoretyczne i~technologie, na których
oparto projekt: model komponentowy i~cykl życia obiektów w~Unity, system fizyki,
wstrzykiwanie zależności, programowanie reaktywne, wzorzec MVVM oraz wzorce projektowe
stosowane w~grach. Rozdział trzeci zawiera koncepcję gry, wymagania funkcjonalne
i~niefunkcjonalne, przypadki użycia oraz projekt architektury. W~rozdziale czwartym opisano
implementację poszczególnych modułów wraz z~fragmentami kodu źródłowego. Rozdział piąty
przedstawia sposób testowania gry, wyniki testów funkcjonalnych i~zmiany wprowadzone na ich
podstawie. Rozdział szósty podsumowuje pracę, ocenia stopień realizacji celów i~wskazuje
kierunki dalszego rozwoju. Pracę uzupełniają bibliografia, spisy rysunków, tabel i~listingów
oraz załączniki zawierające informacje o~kodzie źródłowym, pliki konfiguracyjne
i~instrukcję budowania aplikacji.
